다음 그림에서 $\angle a + \angle b = 90^\circ$, $\angle b + \angle c = 90^\circ$이므로 $\angle a = \angle c$

즉, $\triangle\mathrm{ABD}$와 $\triangle\mathrm{CAE}$에서

$\angle\mathrm{ADB} = \angle\mathrm{CEA} = 90^\circ$, $\overline{\mathrm{AB}} = \overline{\mathrm{CA}}$, $\angle\mathrm{ABD} = \angle\mathrm{CAE}$

$\therefore \triangle\mathrm{ABD} \equiv \triangle\mathrm{CAE}$ (RHA 합동)

즉, $\overline{\mathrm{DA}} = \overline{\mathrm{EC}} = 8\,\mathrm{cm}$, $\overline{\mathrm{AE}} = \overline{\mathrm{BD}} = 10\,\mathrm{cm}$이므로

$\overline{\mathrm{DE}} = \overline{\mathrm{DA}} + \overline{\mathrm{AE}} = 8 + 10 = 18\,(\mathrm{cm})$

이때, 사각형 $\mathrm{DBCE}$는 사다리꼴이므로

(사다리꼴 $\mathrm{DBCE}$의 넓이) $= \dfrac{1}{2}\times (\overline{\mathrm{BD}} + \overline{\mathrm{CE}})\times \overline{\mathrm{DE}}$

$= \dfrac{1}{2}\times (10 + 8)\times 18 = 162\,(\mathrm{cm}^2)$

$\therefore \triangle\mathrm{ABC}$ = (사다리꼴 $\mathrm{DBCE}$의 넓이) $- (\triangle\mathrm{ABD} + \triangle\mathrm{CAE})$

= (사다리꼴 $\mathrm{DBCE}$의 넓이) $- 2\triangle\mathrm{ABD}$

$= 162 - 2\times \dfrac{1}{2}\times 10\times 8 = 82\,(\mathrm{cm}^2)$
